\subsection{Aprendizaje profundo en tareas de visión por computadora}

Las limitaciones que presentan los enfoques tradicionales en problemas de clasificación de imágenes hacen evidente la necesidad de recurrir a métodos con mayor capacidad de representación y aprendizaje. En este sentido, el aprendizaje profundo ha adquirido un papel central en la visión por computadora, al ofrecer modelos capaces de extraer y jerarquizar características directamente a partir de los datos visuales. Por ello, en el presente apartado se introduce el aprendizaje profundo como paradigma de especial relevancia para el análisis de imágenes y para la solución del problema planteado en este trabajo.

\subsubsection{Qué es el aprendizaje profundo}

El aprendizaje profundo es una rama del aprendizaje automático basada en redes neuronales artificiales con múltiples capas, diseñadas para aprender patrones complejos a partir de grandes volúmenes de datos. A diferencia de los enfoques tradicionales, en los que un especialista debía definir manualmente las características visuales relevantes, el aprendizaje profundo permite que el propio modelo descubra automáticamente los rasgos más útiles para resolver una tarea. En visión por computadora, este enfoque ha resultado especialmente valioso porque facilita el aprendizaje de representaciones jerárquicas desde la imagen de entrada hasta la decisión final del sistema \cite{omahony2019deep_learning_vs_traditional_cv}.

Además, el aprendizaje profundo se apoya en arquitecturas como las redes neuronales convolucionales, que procesan la información visual mediante filtros capaces de detectar bordes, texturas, formas y patrones cada vez más complejos. Este esquema de aprendizaje de extremo a extremo ha modificado la forma en que se diseñan los sistemas de visión por computadora, ya que el énfasis ya no está en programar reglas explícitas, sino en entrenar modelos con datos anotados para que aprendan directamente de ellos \cite{omahony2019deep_learning_vs_traditional_cv}.

Desde otra perspectiva, el avance del aprendizaje profundo en visión por computadora también se explica por su integración con estrategias que buscan aprovechar mejor los datos disponibles. En contextos donde etiquetar imágenes es costoso o lento, se han desarrollado variantes como el aprendizaje activo profundo, cuyo propósito es seleccionar las muestras más valiosas para anotarlas y así mejorar el entrenamiento con menos esfuerzo humano. Esto muestra que el aprendizaje profundo no solo depende de arquitecturas potentes, sino también de métodos eficientes para gestionar los datos de entrenamiento \cite{wu2022deep_active_learning}.

\subsubsection{Ventajas del aprendizaje profundo en tareas visuales}

Una de las principales ventajas del aprendizaje profundo en tareas visuales es su \textbf{alta capacidad para alcanzar niveles superiores de precisión} en comparación con los métodos tradicionales. Según la literatura revisada, este enfoque ha impulsado mejoras importantes en tareas como clasificación de imágenes, segmentación semántica, detección de objetos e incluso localización y mapeo simultáneo. Esto se debe a que las redes profundas aprenden automáticamente las características más discriminativas, reduciendo la dependencia de la ingeniería manual de atributos \cite{omahony2019deep_learning_vs_traditional_cv}.

Otra ventaja relevante es su \textbf{flexibilidad}. Los modelos basados en aprendizaje profundo pueden volver a entrenarse con conjuntos de datos específicos para distintos problemas, lo que permite adaptarlos a nuevos contextos con relativa facilidad. En contraste, muchos métodos clásicos de visión por computadora tienden a ser más rígidos y dependientes del dominio de aplicación. De esta manera, el aprendizaje profundo ofrece una mayor capacidad de generalización y adaptación cuando se dispone de datos adecuados para el entrenamiento \cite{omahony2019deep_learning_vs_traditional_cv}.

También destaca su \textbf{utilidad en escenarios complejos} donde existen \textit{grandes volúmenes de información visual}. En la revisión sobre aprendizaje activo profundo se señala que el aprendizaje profundo ha permitido avances notables en visión por computadora gracias al uso de conjuntos de datos anotados a gran escala, que proporcionan supervisión directa durante el entrenamiento. Esto ha favorecido el desarrollo de sistemas más robustos para problemas difíciles, aunque también ha puesto en evidencia la necesidad de optimizar el proceso de etiquetado debido a su alto costo en muchas aplicaciones \cite{wu2022deep_active_learning}.

\subsubsection{Aplicación del aprendizaje profundo en visión por computadora}

El aprendizaje profundo se aplica ampliamente en la visión por computadora para resolver tareas de reconocimiento, análisis e interpretación de imágenes. Entre sus aplicaciones más importantes se encuentran \textit{la clasificación de imágenes, la detección de objetos, y la segmentación semántica}. Estas tareas se han beneficiado del uso de redes convolucionales, las cuales han demostrado un rendimiento sobresaliente frente a enfoques tradicionales \cite{omahony2019deep_learning_vs_traditional_cv}.

A nivel más amplio, la revisión sobre aprendizaje activo profundo destaca que la visión por computadora ha experimentado un avance significativo con la incorporación del aprendizaje profundo en aplicaciones que requieren grandes cantidades de datos etiquetados. Entre los campos mencionados se encuentran la conducción autónoma, el análisis de imágenes médicas, la detección de anomalías y otras tareas donde la precisión del modelo resulta crítica. En estos escenarios, el aprendizaje profundo se ha convertido en una herramienta central para extraer información útil a partir de imágenes, videos, nubes de puntos y otros datos visuales \cite{wu2022deep_active_learning}.

Asimismo, su aplicación no se limita a tareas bidimensionales convencionales. También se ha extendido a contextos emergentes como la \textit{visión 3D, la percepción robótica y los sistemas híbridos} que combinan técnicas tradicionales con redes profundas. Esto confirma que el aprendizaje profundo no solo ha transformado la visión por computadora actual, sino que también sigue ampliando su alcance hacia nuevos problemas y modalidades de datos \cite{omahony2019deep_learning_vs_traditional_cv}.

\subsubsection{Importancia de los datos anotados}

Los datos anotados son fundamentales en el aprendizaje profundo porque proporcionan la supervisión necesaria para que el modelo aprenda a relacionar una entrada visual con una salida correcta. En visión por computadora, las imágenes etiquetadas permiten entrenar sistemas capaces de clasificar, detectar o segmentar con mayor precisión. De acuerdo con las fuentes revisadas, el crecimiento del aprendizaje profundo ha estado fuertemente vinculado a la disponibilidad de grandes conjuntos de datos anotados, ya que estos hacen posible entrenar modelos complejos con un mejor desempeño \cite{omahony2019deep_learning_vs_traditional_cv}, \cite{wu2022deep_active_learning}.

Sin embargo, aunque los datos anotados son indispensables, su obtención representa uno de los principales desafíos del área. En muchas aplicaciones, el etiquetado manual requiere tiempo, conocimiento especializado y altos costos. La revisión sobre aprendizaje activo profundo subraya que en tareas como la segmentación urbana o el análisis de imágenes médicas, la anotación puede ser extremadamente demandante, incluso hasta superar los 90 minutos por imagen en ciertos casos. Esto convierte a los datos anotados en un recurso valioso y limitado, cuya calidad influye directamente en el rendimiento del modelo \cite{wu2022deep_active_learning}.

Por esa razón, no solo importa disponer de datos anotados, sino contar con datos bien seleccionados y de alta calidad. En este sentido, enfoques como el aprendizaje activo profundo buscan reducir la carga de anotación al priorizar las muestras más informativas para el entrenamiento. Así, la importancia de los datos anotados no radica únicamente en su cantidad, sino en su relevancia, precisión y utilidad para mejorar la capacidad de generalización de los modelos de visión por computadora \cite{wu2022deep_active_learning}.